\documentclass[slidetop,11pt]{beamer}
%
% Ces deux lignes à décommenter pour sortir 
% le texte en classe article
% \documentclass[class=article,11pt,a4paper]{beamer}
% \usepackage{beamerbasearticle}

% Packages pour les français
%
\usepackage[T1]{fontenc} 
\usepackage[utf8]{inputenc}
\usepackage[frenchb]{babel}
% pour un pdf lisible  l'cran si on ne dispose pas 
% des fontes cmsuper ou lmodern
\usepackage{lmodern}
%\usepackage{aeguill}

% Pour afficher le pdf en plein ecran
% (comment pour imprimer les transparents et pour les tests)
%\hypersetup{pdfpagemode=FullScreen}

\usepackage{listings}   % need for code encapsulation

\lstset{
	language=Python,
	numbers=left, numberstyle=\tiny, stepnumber=1, numbersep=7pt,
	morekeywords={%
    %%% BOUCLE, TEST & Co.
      if, elif, then, else, and, or, until, yield,repeat,
    %%% IMPORT & Co.
      import, from, package,
    %%% FONCTIONS NUMERIQUES
      cos, sin, tan, acos, asin, atan,
    %%% CONSTANTES
      pi, True, False,
    %%% Types .
      Int, String, Array, Integer, 
    %%% Mots-Clefs supp.
      object, class, protected, public, private, final, val, var, new, ?, !
    },
    sensitive=true,
    morecomment=[l]{\#},
	keywordstyle=\bfseries\emph,
	breaklines=true,	
	frameround=fttt,
	basicstyle= \mdseries\scriptsize }
%\lstset{language=bash,basicstyle=\scriptsize,commentstyle=\small\itshape,stringstyle=\ttfamily,numbers=left,numberstyle=\tiny,stepnumber=1,showstringspaces=false}
	

% ------------------------------------------------
%-----------   styles pour beamer   --------------
% ------------------------------------------------
%
% ------------- Choix des couleurs ---------------
%\xdefinecolor{fondtitre}{rgb}{0.20,0.43,0.09}  % vert fonce
% la mme d'une autre manire
\definecolor{fondtitre}{HTML}{336E17}

%\xdefinecolor{coultitre}{rgb}{0.41,0.05,0.05}  % marron
% la mme d'une autre manire
\definecolor{coultitre}{RGB}{105,13,13}

\xdefinecolor{fondtexte}{rgb}{1,0.95,0.86}     % ivoire

% Redfinit la couleur de fond pour imprimer sur transparents
%\xdefinecolor{fondtexte}{rgb}{1,1,1}     % blanc

% commande differente pour les couleurs nommes - de base
%\colorlet{coultexte}{black} 

% -------------- Fioritures de style -------------
% Fait afficher l'ensemble du frame 
% en peu lisible (gris clair) ds l'ouverture
\beamertemplatetransparentcovered

% Supprimer les icones de navigation (pour les transparents)
%\setbeamertemplate{navigation symbols}{}

% Mettre les icones de navigation en mode vertical (pour projection)
%\setbeamertemplate{navigation symbols}[vertical]

% ------------ Choix des thmes ------------------
%\usecolortheme{crane}
%\usetheme{darmstadt}
%\usetheme{frankfurt}
%\useoutertheme{default}
%\useinnertheme{default}
%\useinnertheme[shadow=true]{rounded}
\mode<presentation>
\useoutertheme{tree}
\usecolortheme{whale}
\usecolortheme{orchid}
\useinnertheme[shadow=true]{rounded}

\setbeamerfont{block title}{size={}}

% Dfinition de boites en couleur spcifiques
% premire mthode
\setbeamercolor{bas}{fg=coultitre, bg=fondtitre!40}
\setbeamercolor{haut}{fg=fondtitre!40, bg=coultitre}
% deuxime mthode
\beamerboxesdeclarecolorscheme{clair}{fondtitre!70}{coultitre!20}
\beamerboxesdeclarecolorscheme{compar}{coultitre!70}{fondtitre!20}

% insrer le nombre de pages
%\logo{\insertframenumber/\inserttotalframenumber}

%------------ fin style beamer -------------------

% Faire apparatre un sommaire avant chaque section
% \AtBeginSection[]{
%   \begin{frame}
%   \frametitle{Plan}
%   \medskip
%   %%% affiche en dbut de chaque section, les noms de sections et
%   %%% noms de sous-sections de la section en cours.
%   \small \tableofcontents[currentsection, hideothersubsections]
%   \end{frame} 
% }

% ----------- Contenu de la page de titre --------
\title{Programmation Concurrente}
\subtitle{Communicating Scala Objects\\Java Communicating Sequential Processes}
%\subtitle{Fichier test}
\author{Jérôme \textsc{Gazel} \\ Clément \textsc{Schiano de Colella}}
\institute{\'Ecole Centrale de Nantes}
\date{\oldstylenums{vendredi 25 mars 2011}}
\logo{\includegraphics[height=1cm]{Logo_ECN.png}}
% ------------------------------------------------
% -------------   Début document   ---------------
% ------------------------------------------------
\begin{document}
%--------- Ecriture de la page de titre ----------
% avec la commande frame simplifie
\frame{\titlepage}
%
\part{Projet d'Application} 
%------------------- Sommaire ----------------
\begin{frame}{Sommaire}
  \small \tableofcontents[hideallsubsections]
\end{frame} 


% ---------------------------------------------
% ------------ INTRODUCTION-------------
% ---------------------------------------------
\section{Introduction}
% Sommaire local. En deux colonnes
\begin{frame}{Plan}
  \tableofcontents[sections=\thesection]
\end{frame}
% ---------------------------------------------
\section{Programmation concurrente}
% avec l'environnement frame
\begin{frame}[label=pagesimple]
  \frametitle{Programmation concurrente}
    \begin{itemize}[<+->]
  \item Qu'est-ce que la programmation concurrente?
  \item Quelles sont ces intérêts?
   \item Quels sont les outils disponibles?
    \end{itemize}
\end{frame}

\subsection{Présentation de CSP}
\begin{frame} 
  \frametitle{Présentation de CSP}
  1978 $\rightarrow$ Article célèbre de Hoare\\ 
  \pause
  \begin{block}{CSP}
  CSP est un langage :
\begin{itemize}[<+->]
\item impératif
\item parallèle
\item communications par échanges synchrones de message
 \end{itemize}
 \end{block}
\end{frame}

\begin{frame} [containsverbatim]
\framesubtitle{Notion de rendez-vous}
Exemple d'échange de message :
\begin{lstlisting}[frame=trBL]
P:: Q ! M(x)
Q:: P ? M(y)
\end{lstlisting}

 \begin{lstlisting}[frame=trBL]
<commande gardee> ::= <garde> -> <suite d instructions>
<garde> ::= <CB>, <CES>
\end{lstlisting}
\end{frame}

\begin{frame} [containsverbatim]
\framesubtitle{Notion de garde et commande gardée}



\begin{lstlisting}[frame=trBL,title={Sélective}]
[[
(state=1) -> state:=state+1;
[]
(state=1) -> state:=state-1;
]]
\end{lstlisting}

\begin{lstlisting}[frame=trBL,title={Répétitive}]
x:=1;
*[[
(x<=N) -> imprimer(x); x++;
[]
(x=N) -> skip;
]]
\end{lstlisting}
\end{frame}

\section{Implémentation des différents langages}
\subsection{Communicating Scala Objects – CSO}
\begin{frame} 
  \frametitle{Implémentation des différents langages}
  \framesubtitle{Communicating Scala Objects – CSO}
  \begin{quotation}
     \textit{\og CSO (Communicating Scala Objects) – a notationally convenient embedding of the essence of occam in a modern, generically typed, object-oriented programming language that is compiled to Java Virtual Machine (JVM) code.\fg}
   \begin{flushright}
\textsc{Bernard Sufrin}
\end{flushright}
\end{quotation}   
\end{frame}

\subsection{Installation de Scala et CSO}
\begin{frame} 
\frametitle{Installation de Scala et CSO}
Très peu de documentation disponible
\pause
\begin{quotation}
\textit{\og This task speaks for itself. If you find it hard to do this then you may find it impossible
to accomplish any of the subsequent tasks.\fg}
\begin{flushright}
\textsc{Bernard Sufrin}
\end{flushright}
\end{quotation}
\end{frame}

\begin{frame} 
\frametitle{Lyfe}
Exemple d'execution
\end{frame}

\begin{frame} 
\frametitle{Premiers essais en CSO...}

\end{frame}

\begin{frame} 
\frametitle{...et premiers problèmes}

\end{frame}

\begin{frame} 
\frametitle{Les différents pipes}
\begin{block}{Canaux synchronisés}
\begin{itemize}[<+->]
\item \textsf{OneOne[T]} 
\item \textsf{ManyOne[T]} 
\item \textsf{OneMany[T]} 
\item \textsf{ManyMany[T]} 
\end{itemize}
\end{block}
\end{frame}

\begin{frame} 
\frametitle{Algorithmes classiques}

\end{frame}

\subsection{Java Communicating Sequential Processes – JCSP}
\begin{frame}
\frametitle{Java Communicating Sequential Processes – JCSP}
\begin{quotation}
\textit{\og Although CSP is a mathematical system, JCSP does not require in-depth mathematical skill, allowing instead that programmers can achieve well-behaved software just by following simple rules.\fg}
\begin{flushright}
\textsc{University of Kent at Canterbury}
\end{flushright}
\end{quotation}
\end{frame}

\subsection{Exemples fournis avec JCSP}
\begin{frame}
\frametitle{Exemples fournis avec JCSP}
Exemple d'execution
\end{frame}

\subsection{Producteur/Consommateur}
\begin{frame}[containsverbatim]
\framesubtitle{Producteur/Consommateur}
\begin{lstlisting}[frame=trBL,title={Producteurs-Consommateurs: Producteur.java}]
import org.jcsp.lang.*;
public class Producteur implements CSProcess
{
  final private ChannelOutput out;
  public Producteur (ChannelOutput out)
        {
                this.out = out ;
        }
        public void run ()
        {
                for (int i=1;i<=100 ; i=i+1)
                {
                        out.write (i);
                }
        }
}
\end{lstlisting}
\end{frame}

\begin{frame}[containsverbatim]
\framesubtitle{Producteur/Consommateur}
\begin{lstlisting}[frame=trBL,title={Producteurs-Consommateurs: Consommateur.java}]
import org.jcsp.lang.*;
public class Consommateur implements CSProcess
{
        final private ChannelInput in;
        public Consommateur (ChannelInput in)
        {
                this.in = in;
        }
        public void run ()
        {
                while (true)
                {
			Integer d = (Integer) in.read();
                        System.out.println ("Lit :" +d);
                }
        }
}
\end{lstlisting}
\end{frame}

\begin{frame}[containsverbatim]
\framesubtitle{Producteur/Consommateur}
\begin{lstlisting}[frame=trBL,title={Producteurs-Consommateurs: Main.java}]
import org.jcsp.lang.*;

public class Main
{
        public static void main (String[] args)
        {
	final One2OneChannel c = Channel.one2one();
             new Parallel
		(
		new CSProcess[]
			{
			new Producteur(c.out()) ,
			new Consommateur(c.in())
			} 
		).run ();
	}
}
\end{lstlisting}
\end{frame}

\subsection{D\^iner des philosophes}
\begin{frame}
Exemple d'execution
\end{frame}

\section{Comparaisons}
\subsection{Comparaisons - CSO}
\begin{frame}
\frametitle{Comparaisons - CSO}
\begin{itemize}[<+->]
\item Avantages de CSO
\begin{itemize}[<+->]
\item CSO est moins verbeux que JCSP. Les programmes sont donc souvent plus courts ...
\item ceci permet une meilleure maintenance du code,
\item une meilleure lisibilité, ce qui facilite le partage
\item enfin l'argument majeur est les possibilités offertes par la programmation fonctionelle
\end{itemize}
\item Inconvénients de CSO
\begin{itemize}[<+->]
\item une documentation restreinte
\item le caractère quasi inexistent d'une communauté active de développement
\item le statut qui nous semble trop exclusif de Bernard Sufrin.
\end{itemize}
\end{itemize}
\end{frame}

\subsection{Comparaisons - JCSP}
\begin{frame}
\frametitle{Comparaisons - JCSP}
\begin{itemize}[<+->]
\item Avantages de JCSP
\begin{itemize}[<+->]
\item une simplicité d'usage m\^eme pour ceux ne disposant pas de connaissances particulières dans les langages de programmation concurrente.
\item une installation bien plus aisée que celle de CSO (cet argument est uniquement valable si l'on ne dispose pas du fichier d'installation pour CSO que nous avons conçu)
\item une documentation riche et fournie
\item de nombreux forums permettant l'entraide et la résolution de problèmes
\end{itemize}
\item Inconvénients de JCSP
\begin{itemize}[<+->]
\item le caractère trop restreint de son utilisation
\item l'incertitude concernant sa pérennité
\end{itemize}
\end{itemize}
\end{frame}


\section{Conclusion}
\begin{frame}
\frametitle{Conclusion}
\begin{itemize}[<+->]
\item \'Etudier en cours, même de façon succinte Scala, alternative de plus en plus crédible de Java,
\item Poursuivre l'étude de la programmation concurrente mais il nous semble encore précoce de traiter de CSO en cours de PASYR,
\item \'Eventuellement proposer un sujet de veille technologique concernant les solutions existantes de programmation concurrente.
\end{itemize}
\end{frame}

\section*{Bibliographie}

\begin{frame} 
\frametitle{Bibliographie}

\begin{thebibliography}{Cpa2008}
\bibitem[Pasyr]{pasyr} {\textit{Cours de Parallélisme et Systèmes Répartis}, Olivier Roux, \'Ecole Centrale de Nantes, 2011}

\bibitem[Cpa2008]{cpa2008-cso} {\textit{Communicating Process Architectures 2008}, Susan Stepney, Fiona Polack, Alistair McEwan, Peter Welch, and Wilson Ifill (Eds.), IOS Press, 2008
\copyright 2008 The authors and IOS Press. All rights reserved.\\
\url{http://citeseerx.ist.psu.edu/viewdoc/download?doi=10.1.1.164.8856&rep=rep1&type=pdf}}
\end{thebibliography}
\end{frame}

\end{document}
